\section{Limitations}
\label{sec:limitations}

\paragraph{Resource and deployment scope.}
The experiments bound selected logical accounts, not the complete physical
working set. The host limit includes packed and prepared payloads, metadata,
and requested worker stack; OS stack rounding, TLS, allocator overhead,
weak-reference backing, and file-cache residency require additional headroom.
Device samples can miss transient peaks and lack exact process attribution.
The three cache points vary entry and byte limits together, so they neither
isolate a replacement-policy effect nor locate a physical capacity frontier.
The complete compact checkpoint may fit in host memory. Other representations,
long contexts, shared-memory devices, and multimodal workloads remain outside
the evaluated profile.

\paragraph{Statistical and measurement scope.}
The timing study contains five paired blocks of short fixed outputs, not
50 independent repetitions. Shared page and NFS caches can carry state across
fresh servers. The descriptive resampling ranges therefore provide neither
guaranteed coverage nor a significance or power result. Serial-window
goodput excludes startup, gaps, and concurrency. Identical instrumentation
controls one source of bias, but its overhead is unmeasured. The profiling
study observes overlap in two separate requests; it does not identify the
cause of the timing difference. Its target-process checks pass, while
wrapper warnings and an unknown dropped-event count leave global trace
completeness unverified. File-call wall time can include cache service and
preemption, and shared NFS counters do not identify physical backend traffic.

\paragraph{Failure and verification scope.}
The component experiment runs real CUDA operations but injects a software
failure or discards proof after valid fences. It does not exercise a driver
failure, lost hardware event, or full-model uncertain completion. The HTTP
experiment tests a client disconnect followed by a new request, not resumed
execution or recovery from a poisoned owner. Neither experiment establishes
the integrated five-class conservation invariant across a serving request.
Final physical release remains unobserved. Unknown custody can reduce
availability indefinitely; the design supplies no general in-place reset,
durable crash recovery, or arbitrary cross-rank atomic commit.

\paragraph{Reproducibility scope.}
The supplementary materials support offline reconstruction of the public
measurements and their exclusions. They do not independently replay private
process authentication or recover omitted executables, native dependencies,
model payloads, and live mapped-memory contents. Integrity hashes bind retained
inputs but are not external attestations. Earlier rejected or incomplete
studies remain separate rather than becoming successful repetitions.

\section{Conclusion}
\label{sec:conclusion}

Bounded storage-backed inference requires a lifetime contract as well as a
cache policy. \system{} connects source-bound planning and stage-local
admission to exact completion, publication, and custody through retirement.
Its bounded next-expert schedule preserves the blocking control's worker,
budgets, arithmetic, and synchronous consumer proofs. In five paired blocks,
lookahead completes the same checked workload with equal logical reads,
zero wasted read bytes, less host demand waiting, and descriptive reductions
of 8.18\% in end-to-end latency and 11.24\% in time to first content, with
10.91\% higher qualified serial-window goodput. The additional experiments
show bounded completion across three coupled cache budgets, limited measured
read--kernel overlap, safe component abort or retained uncertainty, and
fresh-request recovery after HTTP cancellation. These results support the
mechanism within its measured scope while leaving complete physical
accounting, integrated fault conservation, and final release open.
